% ============================== CI-7 ==============================
\reading{CI-7}{Tokenization and Financial Market Inefficiencies}{DEFER}{0}

\begin{cikeybox}[title={The five objectives, in the spine's own words}]
\begin{enumerate}[label=\textbf{\alph*.}]
\item Explain the process of tokenization and describe the fundamental features of tokenized
assets and digital ledgers.
\item Describe frictions and inefficiencies that can arise during different phases of an
asset's lifecycle and explain how the different features of tokenization can both increase and
mitigate risks related to these frictions.
\item Explain how the use of tokenization and digital ledgers can impact financial market
externalities, including those related to the transmission of shocks across the financial
system, investments in market infrastructure, and network and knowledge effects.
\item Explain how the proliferation of distributed ledger technologies can impact retail
investors and brokerage firms.
\item Compare different models and methods that are currently available to achieve sharedness,
programmability, and trust on a ledger, and describe the conditions under which trust can fail
in such models.
\end{enumerate}
\greyline{Source: ``Tokenization and Financial Market Inefficiencies'', IMF Fintech Notes
(January 2025). The class notes label this Reading 106; it is CI-7 on the 2026 spine. Like
CI-6 it carries a DEFER tag and an SRC of 0, and like CI-6 it is built in full.}
\end{cikeybox}

% ------------------------------------------------------------------
\lo{CI-7 a}{Explain the process of tokenization and describe the fundamental features of tokenized assets and digital ledgers.}

\subsection*{Tokens and tokenization}

\begin{itemize}
\item \term{Tokens} are digital representations of ownership recorded on a ledger. Think of
them as a digital certificate that proves who owns an asset and allows easy transfer between
parties.
\item \term{Tokenization} is the process of converting assets into digital tokens, making them
easier to transfer, automate, and manage. It is not a single technology, but a set of tools
built on distributed ledgers.
\end{itemize}

\subsection*{Types of tokens}

\begin{enumerate}
\item \term{Native tokens}. Examples: Bitcoin, Ether, digital art.
\item \term{Non-native tokens}. Examples: equities, bonds, commodities, real estate.
\end{enumerate}

\begin{center}
\renewcommand{\arraystretch}{1.25}
\begin{tabularx}{\textwidth}{@{}>{\raggedright\arraybackslash}X >{\raggedright\arraybackslash}X@{}}
\toprule
\rowcolor{PaleNavy}
{\sffamily\bfseries\color{FRMNavy}Tokenization} &
{\sffamily\bfseries\color{FRMNavy}Token creation}\\
\midrule
\textbf{Creation of assets or representations of assets on a digital ledger that is:}
\begin{itemize}[leftmargin=12pt, topsep=3pt, itemsep=1pt]
\item shared
\item programmable
\item trusted
\end{itemize}
&
\textbf{Can be attained through:}
\begin{itemize}[leftmargin=12pt, topsep=3pt, itemsep=1pt]
\item new issuance (native)
\item new virtual representation of an existing asset (non-native)
\item destruction of an existing asset and reissuance on a ledger (migration)
\end{itemize}\\
\bottomrule
\end{tabularx}
\end{center}
\figcap{Table CI-7.1 (the notes' Figure 2). The left column defines the \emph{state} a ledger
must be in and the right column lists the three \emph{routes} to getting an asset onto it.
Migration is the one most easily forgotten and the one with the most legal weight, since the
original asset is destroyed rather than mirrored. Source: IMF Fintech Notes.}

\subsection*{Key features of tokenized systems}

Tokenization functions effectively only when the underlying ledger provides three essential
features.

\begin{cifigblock}
{\centering
\begin{tikzpicture}[
  font=\sffamily\scriptsize,
  hd/.style={font=\sffamily\bfseries\scriptsize\color{FRMNavy}, align=center},
  def/.style={draw=FRMRule, fill=PaleNavy, text width=3.4cm, align=center,
              minimum height=1.5cm, inner sep=4pt, rounded corners=1pt},
  res/.style={draw=FRMRule, fill=PaleGold, text width=3.4cm, align=center,
              minimum height=1.25cm, inner sep=4pt, rounded corners=1pt}
]
\node[hd] (h1) at (0,0)   {1 \ \textbullet\ \ Trust};
\node[hd] (h2) at (3.9,0) {2 \ \textbullet\ \ Sharedness};
\node[hd] (h3) at (7.8,0) {3 \ \textbullet\ \ Programmability};

\node[def, below=3mm of h1] (d1) {ownership and transactions are recorded accurately and
  cannot be altered without the owner's consent};
\node[def, below=3mm of h2] (d2) {all relevant parties can access the ledger, hold assets
  recorded on it, and transfer ownership};
\node[def, below=3mm of h3] (d3) {smart contracts can automate instructions and embed
  financial logic directly into assets};

\node[res, below=3mm of d1] (o1) {ownership is secure, transactions are accurate};
\node[res, below=3mm of d2] (o2) {accessible to all, direct ownership updates};
\node[res, below=3mm of d3] (o3) {automatic payments, conditional transfers, compliance};

\foreach \a/\b in {h1/d1, h2/d2, h3/d3, d1/o1, d2/o2, d3/o3}
  \draw[-{Stealth[length=4pt]}, FRMGold, line width=0.7pt] (\a) -- (\b);
\end{tikzpicture}\par}
\figcap{Figure CI-7.1. The feature is in the middle row, what it buys you is in the bottom
row. The three are not interchangeable and each answers a different question: trust is about
whether the record can be altered, sharedness is about who can reach it, and programmability
is about whether the ledger can act on its own. A ledger missing any one of the three cannot
support tokenization, which is the claim the objective turns on.}
\end{cifigblock}

% ------------------------------------------------------------------
\lo{CI-7 b}{Describe frictions and inefficiencies that can arise during different phases of an asset's lifecycle and explain how the different features of tokenization can both increase and mitigate risks related to these frictions.}

\subsection*{Market frictions in traditional asset markets}

Market frictions are inefficiencies that arise across an asset's lifecycle, namely issuance,
trading, and servicing. These frictions increase costs, delays, and risks for market
participants.

\begin{cifigblock}
{\centering
\begin{tikzpicture}[
  font=\sffamily\scriptsize,
  qd/.style={draw=FRMRule, fill=PaleNavy, text width=3.0cm, align=left, inner sep=4pt,
             rounded corners=1pt},
  seg/.style={fill=FRMNavy!80, draw=white, line width=1.2pt}
]
\fill[seg] (0,0) -- (0,2.0) arc[start angle=90, end angle=180, radius=2.0] -- cycle;
\fill[seg] (0,0) -- (2.0,0) arc[start angle=0,  end angle=90,  radius=2.0] -- cycle;
\fill[seg] (0,0) -- (0,-2.0) arc[start angle=270, end angle=360, radius=2.0] -- cycle;
\fill[seg] (0,0) -- (-2.0,0) arc[start angle=180, end angle=270, radius=2.0] -- cycle;
\node[white, font=\sffamily\bfseries\tiny] at (-0.88,0.92) {Issuance};
\node[white, font=\sffamily\bfseries\tiny] at ( 0.88,0.92) {Exchange};
\node[white, font=\sffamily\bfseries\tiny] at ( 0.88,-0.92) {Servicing};
\node[white, font=\sffamily\bfseries\tiny] at (-0.88,-0.92) {Redemption};

\node[qd] (q1) at (-4.45,1.55)
  {\textbf{Issuance}\\ asymmetric information \textbullet\ search friction \textbullet\
   transaction costs};
\node[qd] (q2) at ( 4.45,1.55)
  {\textbf{Exchange}\\ search friction \textbullet\ transaction costs \textbullet\
   counterparty risk};
\node[qd] (q3) at ( 4.45,-1.55) {\textbf{Servicing}\\ transaction costs};
\node[qd] (q4) at (-4.45,-1.55) {\textbf{Redemption}\\ transaction costs};
\draw[FRMRule, line width=0.7pt] (q1.east) -- (-1.45,1.38);
\draw[FRMRule, line width=0.7pt] (q2.west) -- ( 1.45,1.38);
\draw[FRMRule, line width=0.7pt] (q3.west) -- ( 1.45,-1.38);
\draw[FRMRule, line width=0.7pt] (q4.east) -- (-1.45,-1.38);
\end{tikzpicture}\par}
\figcap{Figure CI-7.2 (the notes' Figure 3). Frictions along the lifecycle of an asset. Note
that the load is not evenly spread: issuance and exchange carry three frictions each, while
servicing and redemption carry only transaction costs. Counterparty risk appears at exactly
one point, exchange, which is why it is the friction tokenization cannot fully remove.}
\end{cifigblock}

\begin{center}
\renewcommand{\arraystretch}{1.25}
\begin{tabularx}{\textwidth}{@{}p{3.0cm} >{\raggedright\arraybackslash}X@{}}
\toprule
\rowcolor{FRMSoft}
{\sffamily\bfseries\color{FRMNavy}Key friction} &
{\sffamily\bfseries\color{FRMNavy}What it is}\\
\midrule
\term{Information asymmetry} & One party has more information (e.g., IPOs)\\
\term{Search frictions} & Difficulty finding buyers and sellers\\
\term{Transaction costs} & Fees, time, and execution costs\\
\term{Counterparty risk} & Risk that the other party defaults\\
\bottomrule
\end{tabularx}
\end{center}

\subsection*{Role of intermediaries}

Specialised institutions reduce these frictions but add costs.

\begin{enumerate}[label=\textbf{\alph*)}]
\item \term{Investment banks}: reduce information asymmetry through due diligence and
underwriting.
\item \term{CSDs and clearinghouses}: manage settlement, reduce counterparty risk.
\item \term{CCPs}: act as intermediary in trades, absorb default risk.
\end{enumerate}

\subsection*{Externalities, internalities, and market power}

Transacting in financial markets also gives rise to externalities, internalities, and market
power issues.

\begin{enumerate}
\item \term{Externalities}.
  \begin{itemize}
  \item Positive: network effects, so more participants means more liquidity.
  \item Negative: moral hazard, so risk-taking due to expected bailouts.
  \end{itemize}
\item \term{Internalities}. Poor decisions by investors, for example trading without
understanding risk.
\item \term{Market power}. Firms charging higher fees due to switching barriers.
\end{enumerate}

\subsection*{Mitigating market frictions through tokenization}

\begin{enumerate}
\item \term{Asset issuance.}
  \begin{enumerate}[label=\alph*)]
  \item Traditional issuance relies on intermediaries like registrars, banks, and trust
  companies for recordkeeping and transactions, increasing costs and complexity.
  \item Tokenization links assets directly to owners on a shared ledger, reducing
  recordkeeping and custodial needs.
  \item Smart contracts automate processes like share auctions, saving time and resources.
  \item Key data can also be embedded within tokens, improving transparency and investor
  access.
  \end{enumerate}
\item \term{Asset trading.}
  \begin{enumerate}[label=\alph*)]
  \item Traditional trading involves counterparty risk, delayed settlement, and transaction
  costs, even with institutions like CSDs and clearinghouses.
  \item Tokenization allows simultaneous or near-instant settlement through smart contracts. A
  trade is executed only when both cash and securities are available, reducing settlement
  delays, lowering costs, and improving efficiency. This is especially useful in complex
  transactions such as cross-currency trades.
  \item A shared ledger also reduces search frictions, particularly in OTC markets, by making
  it easier for participants to find counterparties and improving transparency. This can narrow
  bid-ask spreads and reduce dependence on broker-dealers.
  \item However, tokenization does \emph{not} completely eliminate counterparty risk before
  contract execution, which may still require CCPs or collateral, which introduces opportunity
  costs.
  \end{enumerate}
\item \term{Asset servicing and redemption.}
  \begin{enumerate}[label=\alph*)]
  \item In traditional markets, servicing assets involves costs related to paying interest,
  principal, and dividends, along with maintaining accurate ownership records.
  \item Tokenization can automate these functions through smart contracts. Payments can be
  programmed based on token ownership, payment date, and amount due. This reduces
  administrative burden, lowers servicing costs, and improves operational efficiency.
  \end{enumerate}
\end{enumerate}

\begin{citrapbox}[title={The objective says ``both increase and mitigate''}]
Three of the four items above are mitigation. The one that runs the other way is asset
trading item (d): counterparty risk before contract execution survives tokenization, so CCPs
or collateral may still be needed, and the collateral itself carries an opportunity cost. The
fuller set of \emph{increased} risks sits under objective (c).
\end{citrapbox}

% ------------------------------------------------------------------
\lo{CI-7 c}{Explain how the use of tokenization and digital ledgers can impact financial market externalities, including those related to the transmission of shocks across the financial system, investments in market infrastructure, and network and knowledge effects.}

Tokenization improves \term{liquidity, speed, and interconnectedness}, but also introduces
\term{systemic risks} through externalities, leverage, and rapid transmission of shocks.

\subsection*{1. Shock transmission externalities}

In traditional markets, leverage (high debt), interconnectedness (a shock affecting one
institution can trigger a domino effect) and liquidity mismatches can amplify financial
shocks. \term{Tokenization can intensify this effect.}

\begin{enumerate}[label=\textbf{\alph*)}]
\item \term{Faster transactions}: quicker spread of shocks.
\item \term{Lower issuance cost of tokens}: encourages higher leverage, hence increases
default and systemic risk.
\item \term{Retail funds shift away from bank deposits}: leads banks to rely on unstable
wholesale funding.
\item \term{Complex assets and rehypothecation} (using the same assets as collateral): more
interconnected risk chains.
\item Faster transaction speeds reduce the time available for policymakers to intervene,
allowing shocks to spread more rapidly.
\item Automated trading via smart contracts can trigger rapid, large-scale transactions,
increasing the likelihood of flash crashes and reinforcing downward price spirals.
\end{enumerate}

\subsection*{2. Market infrastructure investment externalities}

\begin{enumerate}[label=\textbf{\alph*)}]
\item As reliance on digital ledgers grows, so does the importance of robust infrastructure
and cybersecurity.
\item Underinvestment in security by private operators can create system-wide vulnerabilities
affecting all participants.
\item Risks include system failures, fraud, human and software errors, and cyberattacks.
\item Concentration risk arises when a few dominant ledgers become critical points of failure,
increasing systemic exposure.
\end{enumerate}

\subsection*{3. Network externalities}

\begin{enumerate}[label=\textbf{\alph*)}]
\item More users bring greater liquidity, better price discovery, lower transaction costs, and
narrower bid-ask spreads.
\item Tokenization enables fractional ownership across a wide range of assets, improving retail
participation and diversification.
\item However, if multiple non-interoperable ledgers emerge, markets can become fragmented into
isolated liquidity pools.
\item Such fragmentation increases search frictions, reduces efficiency, and weakens the
overall benefits of tokenization.
\end{enumerate}

\subsection*{4. Knowledge externalities}

\begin{center}
\renewcommand{\arraystretch}{1.25}
\begin{tabularx}{\textwidth}{@{}p{2.0cm} >{\raggedright\arraybackslash}X@{}}
\toprule
\rowcolor{FRMSoft}
&{\sffamily\bfseries\color{FRMNavy}Knowledge externalities}\\
\midrule
\textcolor{FRMGreen}{\bfseries Positive} & Innovators can build upon existing smart contracts
at low cost, accelerating technological progress for the whole community\\
\textcolor{FRMRed}{\bfseries Negative} & Because code can be easily copied, individual firms
may have less incentive to invest in expensive R\&D, as they cannot easily protect their
intellectual property without using restrictive patents\\
\bottomrule
\end{tabularx}
\end{center}
\figcap{Table CI-7.2. The same property, that code on a shared ledger is visible and copyable,
produces both rows. Faster diffusion and weaker incentive to fund the original research are
not two separate effects to weigh against each other; they are one mechanism seen from the
copier's side and the inventor's side.}

% ------------------------------------------------------------------
\lo{CI-7 d}{Explain how the proliferation of distributed ledger technologies can impact retail investors and brokerage firms.}

Distributed ledger technology (DLT) changes how assets are \term{traded, transferred, and
serviced}, affecting both retail investors and brokerage firms.

\subsection*{1. Retail investor internalities}

Direct access to trading on tokenized ledgers can lead retail investors to make suboptimal
decisions. By bypassing financial advisors to save fees, investors may invest in complex or
risky assets they do not fully understand. This can result in poor diversification and higher
exposure to sudden losses, especially in the absence of strong investor protection frameworks.

\subsection*{2. Market power inefficiencies}

\begin{itemize}
\item Market power inefficiencies arise when limited competition allows dominant players to
charge higher fees.
\item In traditional markets, brokers retain this power because switching is cumbersome.
Processes like in-kind transfers require intermediaries such as clearinghouses, making it slow
and costly for investors to move, so clients remain locked in and fees stay high.
\item DLT reduces these frictions. Instant transfers, transparent ownership, and removal of
intermediaries increase competition, pushing fees lower. Direct trading between investors and
issuers can further pressure brokers, although regulations still require intermediaries in
most markets.
\item However, private ledgers can recreate the same issue. A dominant platform may initially
offer low or free access, but once it gains scale, it can charge monopoly rents and monetise
user data, reintroducing market power inefficiencies.
\end{itemize}

\begin{citrapbox}[title={The argument turns around twice}]
DLT lowers fees by removing the switching barrier, and then a dominant private ledger rebuilds
the barrier at a new location. Any statement that ends at the first turn, ``DLT reduces market
power'', is only half of what the reading says. The same applies to the retail investor: the
benefit is cheaper direct access and the cost is that the advisor who was screening the risk
has been removed along with the fee.
\end{citrapbox}

% ------------------------------------------------------------------
\lo{CI-7 e}{Compare different models and methods that are currently available to achieve sharedness, programmability, and trust on a ledger, and describe the conditions under which trust can fail in such models.}

\subsection*{Tokenization models (sharedness)}

\begin{cifigblock}
{\centering
\begin{tikzpicture}[
  font=\sffamily\scriptsize,
  md/.style={draw=FRMRule, text width=4.8cm, align=left, inner sep=5pt,
             minimum height=1.9cm, rounded corners=1pt}
]
\node[md, fill=PaleNavy] (m1) at (0,0)
  {\textbf{Single ledger model}
   \begin{itemize}[leftmargin=10pt, topsep=2pt, itemsep=1pt]
   \item all participants operate on a single system where assets are issued natively
   \item high efficiency but centralised structure
   \end{itemize}};
\node[md, fill=PaleGold, right=5mm of m1] (m2)
  {\textbf{Common ledger model}
   \begin{itemize}[leftmargin=10pt, topsep=2pt, itemsep=1pt]
   \item assets from other ledgers are represented via escrow certificates
   \item original asset locked, tokenized version traded elsewhere
   \end{itemize}};
\node[md, fill=PaleGreen, right=5mm of m2] (m3)
  {\textbf{Compatible ledger model}
   \begin{itemize}[leftmargin=10pt, topsep=2pt, itemsep=1pt]
   \item different ledgers use standardised technology to interact
   \item enables interoperability across systems without centralisation
   \end{itemize}};
\node[below=3mm of m1.south, font=\sffamily\scriptsize\itshape\color{FRMGrey},
      text width=4.8cm, align=center] {one ledger};
\node[below=3mm of m2.south, font=\sffamily\scriptsize\itshape\color{FRMGrey},
      text width=4.8cm, align=center] {many ledgers, one meeting point};
\node[below=3mm of m3.south, font=\sffamily\scriptsize\itshape\color{FRMGrey},
      text width=4.8cm, align=center] {many ledgers, a common standard};
\end{tikzpicture}\par}
\figcap{Figure CI-7.3. The three models are a spectrum in how they buy sharedness. The single
ledger buys it with centralisation, the common ledger buys it by locking the original asset
and trading a certificate, and the compatible ledger buys it with a standard and no central
point at all. The middle model is the one that introduces a second asset, the escrow
certificate, and therefore a second thing that can fail.}
\end{cifigblock}

\subsection*{Programmability}

Programmability refers to the integration of code and database within the ledger itself. Smart
contracts automate execution of instructions, so there is no need for external systems to
process transactions. This ensures consistency and efficiency in execution.

\subsection*{Trust and trust failures}

\begin{center}
\renewcommand{\arraystretch}{1.25}
\small
\begin{tabularx}{\textwidth}{@{}p{2.9cm} >{\raggedright\arraybackslash}X
                              >{\raggedright\arraybackslash}X@{}}
\toprule
\rowcolor{FRMSoft}
{\sffamily\bfseries\color{FRMNavy}Model} &
{\sffamily\bfseries\color{FRMNavy}How trust is established} &
{\sffamily\bfseries\color{FRMRed}Key trust failures and risks}\\
\midrule
\term{Single-operator system} &
Centralised structure where one entity manages records and validates transactions. Trust
depends on operator credibility and regulatory oversight &
Entire system depends on one entity \textbullet\ high exposure to operational failure or
misuse \textbullet\ no consensus mechanism to distribute trust\\
\addlinespace
\term{Permissionless DLT} &
Decentralised system where anyone can participate in validation. Trust is built through
consensus mechanisms and economic incentives &
51\% attack (majority control of network) \textbullet\ transaction manipulation (front running,
sandwich attacks) \textbullet\ higher exposure to cyber risks, and difficult regulatory
enforcement\\
\addlinespace
\term{Permissioned DLT} &
Hybrid model where only authorised participants validate transactions. Trust is based on
controlled access and oversight by a central authority or consortium &
Risk of collusion among participants \textbullet\ potential manipulation of ledger or rules
\textbullet\ risks mitigated by reputation concerns and regulation\\
\bottomrule
\end{tabularx}
\end{center}
\figcap{Table CI-7.3. Read the middle column first: each model rests trust on a different
thing, and each failure mode in the right column is the failure of exactly that thing. The
single operator rests on one entity's credibility, so the failure is that entity. Permissionless
DLT rests on consensus, so the failure is capturing the consensus. Permissioned DLT rests on a
closed group, so the failure is that group colluding. Note the last row is the only one where
the notes attach a \emph{mitigant} to the risks rather than leaving them open.}
